\medskip

On considère la fonction $f$ définie sur $\R$ par 

\[f(x) = \ln \left(\e^{2x}  - \e^x  + 1\right).\]

On note $\mathcal{C}_f$ sa courbe représentative représentée ci-dessous.

\begin{center}
\psset{unit=1.35cm,arrowsize=2pt 3,comma=true}
\begin{pspicture*}(-3.6,-2)(3,3)
\psaxes[linewidth=1.25pt,Dx=0.5,Dy=0.5,labelFontSize=\scriptstyle]{->}(0,0)(-3.5,-1.95)(3,3)
\psgrid[gridlabels=0pt,subgriddiv=2]
\psplot[plotpoints=2000,linewidth=1.25pt,linecolor=blue]{-3.5}{2}{2.71828 x 2 mul exp 2.71828 x exp sub 1 add ln} \uput[dr](1.4,2.5){\blue $\mathcal{C}_f$}
\end{pspicture*}
\end{center}

Un élève formule les conjectures suivantes à partir de cette représentation graphique :

\smallskip

\begin{tabularx}{\linewidth}{|X|}\hline
\textbf{1.~~}L'équation $f(x) = 2$ semble admettre au moins une solution.\\
\textbf{2.~~}Le plus grand intervalle sur lequel la fonction $f$ semble être croissante est $[- 0,5~;~+\infty[$.\\
\textbf{3.~~}L'équation de la tangente au point d'abscisse $x = 0$ semble être: $y = 1,5x$.\\ \hline
\end{tabularx}

\medskip

Le but de cet exercice est de valider ou rejeter les conjectures concernant la fonction $f$.

\medskip

\textbf{Partie A : Étude d'une fonction auxiliaire}

\medskip

On définit sur $\R$ la fonction $g$ définie par 
\[g(x) = \e^{2x} - \e^x + 1.\]

\begin{enumerate}
\item Déterminer $\displaystyle\lim_{x \to - \infty} g(x)$.
\item Montrer que $\displaystyle\lim_{x \to + \infty} g(x) = + \infty$.
\item Montrer que $g'(x) = \e^x\left(2\e^x - 1\right)$ pour tout $x \in  \R$.
\item Étudier le sens de variation de la fonction $g$ sur $\R$.

Dresser le tableau des variations de la fonction $g$ en y faisant figurer la valeur exacte des extremums s'il y en a, ainsi que les limites de $g$ en $-\infty$ et $+ \infty$.
\item En déduire le signe de $g$ sur $\R$.
\item Sans en mener nécessairement les calculs, expliquer comment on pourrait établir le résultat
de la question 5 en posant $X = \e^x$.
\end{enumerate}

\medskip

\textbf{Partie B}

\medskip

\begin{enumerate}
\item Justifier que la fonction $f$ est bien définie sur $\R$.
\item La fonction dérivée de la fonction $f$ est notée $f'$.

Justifier que $f'(x) = \dfrac{g'(x)}{g(x)}$ pour tout $x \in \R$.
\item Déterminer une équation de la tangente à la courbe au point d'abscisse $0$.
\item Montrer que la fonction $f$ est strictement croissante sur $[-\ln (2)~;~ +\infty[$.
\item Montrer que l'équation $f(x) = 2$ admet une unique solution $\alpha$ sur $[-\ln (2)~;~ +\infty[$ et déterminer une valeur approchée de $\alpha$ à $10^{-2}$ près.
\end{enumerate}

\medskip

\textbf{Partie C}

\medskip

À l'aide des résultats de la partie B, indiquer, pour chaque conjecture de l'élève, si elle est vraie ou fausse. Justifier.

\medskip
